% Mohegan-Pequot (Fielding 2013) — Source Workflow
% Issue #1364
% Status: Research complete, implementation pending

\section{Mohegan-Pequot: Fielding 2013}
\label{sec:1364}

\subsection{Source Description}

\textbf{Recorder:} Stephanie Fielding (Cits Pátunáhshô), Mohegan tribal
member and linguist, with David Costa (Myaamia Center, Miami University of
Ohio).
\textbf{Source:} \emph{Modern Mohegan Dictionary} (2006, revised 2013).
\textbf{Language:} Mohegan-Pequot [xpq] (Y-dialect of SNEA).
\textbf{Corpus:} 145 records (FF-marked entries from Fidelia Fielding's
diaries).
\textbf{Orthography:} Modern practical orthography --- 12 consonant + 6 vowel
symbols, 1:1 phoneme-grapheme correspondence, circumflex for length.

\subsubsection{Fidelia A. H. Fielding (1827--1908)}

Fidelia A. Hoscott Fielding (Mohegan name: \textbf{Jits Bodunaxa} ``Flying
Bird'') was the last fluent speaker of the Mohegan-Pequot language. Born in
1827 at Mohegan, near Norwich, Connecticut, she learned Mohegan as a child
from her grandmother, \textbf{Martha Uncas Shantup} (1761--1859), a direct
descendant of Sachem Uncas. As the number of fluent speakers dwindled through
the late 19th century, Fielding maintained her language by speaking with her
sister and, later, by talking to herself (Prince \& Speck 1904:18--19).

\subsubsection{The Fielding Diaries}

Between 1902 and 1904, Fielding kept a series of diaries written in the
Mohegan language. These are the only known surviving documents produced by a
native speaker of Mohegan-Pequot. The corpus consists of:

\begin{itemize}
  \item \textbf{Three diaries} and a copy of the Lord's Prayer in Mohegan
    (Cornell University Library, acquired 2004 from the Huntington Free
    Library)
  \item A fourth diary was lost in a fire at Columbia University around 1906
    (Granberry 2003:10)
  \item A sermon text published by Prince and Speck (1903) in \emph{American
    Anthropologist}, vol. 5, pp. 193--212
  \item A long narrative tale recorded in phonetic notation by Frank Speck
    (Speck 1904)
\end{itemize}

The diaries were written in Fielding's own orthography --- an English-based
ad-hoc spelling system reflecting her pronunciation. After her death in 1908,
the diaries passed to Frank Speck at the University of Pennsylvania, and later
to Cornell University. In November 2020, Cornell repatriated the diaries to
the Mohegan Tribe in a formal ceremony, where they now reside in the Mohegan
Archives at Uncasville, Connecticut.

\subsubsection{Corpus Structure}

The Brothertown corpus under analysis contains \textbf{145 records} from the
Fielding dictionary. This is a subset of the full dictionary (which numbers
several thousand entries, many reconstructed via comparative Algonquian
methods by David Costa). Records marked \textbf{FF} in the dictionary
indicate words attested directly in Fidelia Fielding's diaries. The 145
records constitute our working corpus because they have explicit
Mohegan-to-phoneme field mappings, whereas the larger set has modern
orthography forms requiring parsing.

\subsubsection{Comparative Significance}

Prince and Speck (1904:18) explicitly compare Fielding to \textbf{Dorothy
Pentreath}, the last speaker of Cornish, noting the ``strangely pathetic''
nature of a dying language spoken only by a single elderly woman who kept it
alive by talking to herself. Despite the ``last stages of decay'' they
observed, they found that ``enough of the original phonetics and grammatical
phenomena'' remained to judge the language's structure.

\subsubsection{Stephanie Fielding's Grammatical Analysis}

Stephanie Fielding (Cits Pátunáhshô), a full-blooded Mohegan tribal member
and descendant of Fidelia Fielding by marriage (her great-great-grandfather's
aunt), served as the Mohegan Tribe's linguist and authored the \emph{Modern
Mohegan Dictionary} (2006, revised 2013). The dictionary was reviewed by David
Costa who researched the lexicon and grammar, developed the Mohegan alphabet,
translated scripts and wrote sample sentences, and proofed the grammar
paradigms.

The dictionary includes:
\begin{enumerate}
  \item \textbf{Pronunciation Guide} --- the practical orthography with
    phoneme mappings
  \item \textbf{Mohegan Grammar Paradigms} --- a systematic grammatical
    sketch covering pronouns, nouns, verbs, and particles
  \item \textbf{Mohegan to English Dictionary} --- stems with part-of-speech
    codes and inflections
  \item \textbf{English to Mohegan Word Finder} --- English lookup with
    Mohegan stems
\end{enumerate}

Fielding's grammar covers the major Algonquian grammatical categories:
pronouns (prefixes and suffixes attached to verb stems, pro-drop/polysynthetic),
nouns (two genders: Animate vs. Inanimate, with plural, obviative, locative,
and absentative inflections), verbs (four orders: Independent, Conjunct,
Imperative, Participial, with transitivity marking VAI, VII, VTA, VTI), and
derivation (person-marking prefixes \textit{nu-} 1st person, \textit{ku-}
2nd person).

\subsubsection{Fielding's Methodology}

Stephanie Fielding worked extensively from texts, particularly the
\textbf{Eliot Bible} (John Eliot's 1663 Massachusett translation), using a
concordance to find Wôpanâak cognates and applying Mohegan-specific sound
changes to reconstruct Mohegan forms. This comparative method is standard for
Algonquian language reclamation. She also drew from Fidelia Fielding's diaries
(primary source, marked FF), sister languages (Wôpanâak, Narragansett, Munsee
Delaware), Proto-Algonquian reconstructions with sound-law application, and
David Costa's grammatical charts and comparative research.

Stephanie Fielding spent the 2017--2018 academic year as a presidential
visiting fellow at Yale University, teaching a course comparing the Mohegan
and Delaware languages. She has also worked with the Unkechaug people of Long
Island to help reclaim their nearly identical language (the primary difference
being the Mohegan /y/ vs. Unkechaug /r/ correspondence).

\subsection{Vowel Inventory}

\subsubsection{Vowel System}

The Mohegan vowel system in the Fielding practical orthography consists of
six plain vowels, four long vowels, and three diphthongs:

\begin{table}[h]
\centering
\begin{tabular}{llll}
\toprule
Grapheme & IPA & Length & Notes \\
\midrule
a & /a/ & short & Low central \\
â & /aː/ & long & \\
e & /iː/ ~ /e/ & short & Front, historically /iː/ lowered; can be /e/ \\
ê & /eː/ & long & Rare; mid front \\
i & /i/ & short & High front \\
o & (not used as plain vowel) & --- & Only occurs in diphthongs \\
ô & /oː/ & long & Very common; mid back rounded \\
u & /u/ & short & High back rounded \\
 & /ə/ & short & Schwa, not written consistently \\
\bottomrule
\end{tabular}
\caption{Mohegan vowel system}
\label{tab:1364-vowels}
\end{table}

\subsubsection{Diphthongs}

\begin{table}[h]
\centering
\begin{tabular}{lll}
\toprule
Grapheme & IPA & Notes \\
\midrule
ay & /aj/ & \\
aw & /aw/ & \\
uy & /uj/ & \\
\bottomrule
\end{tabular}
\caption{Mohegan diphthongs}
\label{tab:1364-diphthongs}
\end{table}

Granberry (2003:14) also reconstructs /oi/ /ɔj/ as marginal diphthongs (as in
English ``boy'') from certain sources.

\subsection{Consonant Inventory}

\begin{table}[h]
\centering
\begin{tabular}{llllll}
\toprule
Grapheme & IPA & Voicing & Place & Manner & Notes \\
\midrule
p & /p/ & Voiceless & Bilabial & Stop & Unaspirated lenis \\
t & /t/ & Voiceless & Dental & Stop & Unaspirated lenis \\
k & /k/ & Voiceless & Velar & Stop & Unaspirated lenis \\
c & /tʃ/ & Voiceless & Palatal & Affricate & Always /tʃ/, never /k/ \\
s & /s/ & Voiceless & Alveolar & Fricative & Fortis/lenis not marked \\
sh & /ʃ/ & Voiceless & Palatal & Fricative & \\
h & /h/ & Voiceless & Glottal & Fricative & Explicitly written \\
m & /m/ & Voiced & Bilabial & Nasal & \\
n & /n/ & Voiced & Alveolar & Nasal & \\
w & /w/ & Voiced & Labiovelar & Approximant & \\
y & /j/ & Voiced & Palatal & Approximant & \\
kw & /kʷ/ & Voiceless & Labiovelar & Stop & Labialized \\
q & /kʷ/ & Voiceless & Labiovelar & Stop & Word-final /kʷ/ \\
' & /ʔ/ & --- & Glottal & Stop & Okina \\
\bottomrule
\end{tabular}
\caption{Mohegan consonant inventory}
\label{tab:1364-consonants}
\end{table}

\subsubsection{Allophonic Variation (from Granberry 2003)}

Granberry's phonological reconstruction (using the ``reconstituted phonetic
units'' methodology of Mary Haas 1954) identifies allophonic patterns:

\begin{itemize}
  \item Stops /p, t, k/ are \textbf{lenis-fortis pairs} overlapping in
    distribution --- the practical orthography uses voiceless symbols but
    actual realization varies
  \item /kw/ word-finally often reduced to [k] before consonants
  \item Vowels show positional length variation: pre-stress short, stressed
    medium-long, post-stress reduced
  \item The phoneme /e/ (plain \textit{e}) ranges from [i] to [e] depending
    on environment, likely reflecting the historical merger of PA *i and *e
    in initial syllables (Oxford 2015)
\end{itemize}

\subsection{Historical Phonology: Proto-Algonquian Reflexes}

\subsubsection{PA Vowel System}

Proto-Algonquian is reconstructed with an eight-vowel system: four qualities
with a length contrast (Bloomfield 1946).

\begin{table}[h]
\centering
\begin{tabular}{lll}
\toprule
PA & Mohegan Reflex & Example \\
\midrule
*a & a & \\
*aː & â & \\
*e & i, a, ə (variable) & \\
*eː & ê & \\
*i & i & \\
*iː & e & \\
*o & u & \\
*oː & ô & \\
\bottomrule
\end{tabular}
\caption{PA to Mohegan vowel reflexes}
\label{tab:1364-pa-vowels}
\end{table}

\subsubsection{The Y-Dialect Classification}

Mohegan-Pequot is classified as a \textbf{Y-dialect} of Southern New England
Algonquian (SNEA). In the traditional classification based on the reflex of
PA *r (sometimes reconstructed as *l):

\begin{table}[h]
\centering
\begin{tabular}{lll}
\toprule
Dialect Type & Reflex of PA *r & Example Language \\
\midrule
N-dialect & /n/ & Massachusett (Wôpanâak), Narragansett, Nipmuc \\
Y-dialect & /y/ & \textbf{Mohegan-Pequot}, Montauk \\
R-dialect & /r/ & Western Connecticut (Loup area) \\
L-dialect & /l/ & Munsee Delaware, Unami \\
\bottomrule
\end{tabular}
\caption{SNEA dialect classification}
\label{tab:1364-dialects}
\end{table}

This means that where Wôpanâak has \textit{n}, Mohegan typically has
\textit{y}. Stephanie Fielding noted this correspondence operating even in
modern Connecticut English pronunciation (e.g., ``Sullivan'' →
``Suh-yuh-van''). Granberry (2003:13) explicitly documents the core
correspondences: Pequot \textit{p, t, k, ch, s} → Mohegan \textit{b, d, g,
j, z} form-initially, and final Pequot \textit{p, t, k, ch, s} → Mohegan
form-final \textit{b, d, g, j, z}.

\subsubsection{Consonant Reflexes}

\begin{table}[h]
\centering
\begin{tabular}{lll}
\toprule
PA & Mohegan & Notes \\
\midrule
*p & p & Unaspirated \\
*t & t & Unaspirated; palatalized to c /tʃ/ before *i, *iː, *y \\
*č & c /tʃ/ & \\
*k & k & Unaspirated \\
*m & m & \\
*n & n & \\
*θ & n (N-dialects), y (Y-dialects) & Key PA distinction \\
*r (controversial) & y & The Y-dialect defining reflex \\
*s & s & \\
*š & sh /ʃ/ & \\
*h & h & Explicit in modern orthography \\
*w & w & \\
\bottomrule
\end{tabular}
\caption{PA to Mohegan consonant reflexes}
\label{tab:1364-pa-consonants}
\end{table}

\subsubsection{Cluster Reflexes}

\begin{table}[h]
\centering
\begin{tabular}{lll}
\toprule
PA Cluster & Mohegan Reflex & Notes \\
\midrule
*ht & ht & Retained in some forms, simplified in others \\
*hk & k & \\
*hp & p & \\
*šk & shk /ʃk/ & \\
*θp & p & With devoicing \\
*mp & p & Nasal cluster simplification \\
*nt & t & \\
*nk & k & \\
*nč & c /tʃ/ & \\
\bottomrule
\end{tabular}
\caption{PA cluster reflexes in Mohegan}
\label{tab:1364-clusters}
\end{table}

\subsection{Comparison with Prince \& Speck (1904)}

\subsubsection{The Prince-Speck Glossary}

J. Dyneley Prince (Columbia University) and Frank G. Speck published
``Glossary of the Mohegan-Pequot Language'' in \emph{American Anthropologist}
(1904, vol. 6, no. 1, pp. 18--44). This glossary contains \textbf{446 words
and forms} collected by Speck from Fidelia Fielding between 1902 and 1903,
submitted to Prince in Fielding's own spelling.

Prince and Speck deliberately preserved Fielding's ``imperfect''
English-orthography system for ``sentimental reasons'' (1904:19), while
providing phonetic transcriptions in parentheses.

\subsubsection{Orthographic Differences}

\begin{table}[h]
\centering
\begin{tabular}{lll}
\toprule
Feature & Fielding (2006) & Prince-Speck (1904) \\
\midrule
Affricate /tʃ/ & \textit{c} always & \textit{ch} or contextual \\
Velar /k/ & \textit{k} & \textit{g} initially, \textit{k, c} elsewhere \\
Voicing & Unaspirated voiceless & Inconsistent (g/k, b/p variation) \\
Vowel length & Circumflex (â, ê, ô) & Colon or context \\
Glottal stop & \texttt{'} (okina) & Inverted comma \\
Schwa & Not written & \texttt{'} (apostrophe) \\
/ʃ/ & \textit{sh} & \textit{sh} \\
Stops & p, t, k (unaspirated) & b/g/p/t/k/d (inconsistent) \\
\bottomrule
\end{tabular}
\caption{Orthographic differences between Fielding and Prince-Speck}
\label{tab:1364-orth-diff}
\end{table}

\subsubsection{The Brothertown Material}

A subset of Prince-Speck's glossary consists of ``Brothertown words'' ---
items collected from an old Mohegan man who had lived at Brothertown, near
Green Bay, Wisconsin, where relocated New England tribes (Tunxis, Wampanoag,
Mohegans, Montauks) had settled. Prince and Speck note these words are
``merely corruptions of Ojibwe forms'' (1904:19) --- relevant since our
corpus bears the name ``Brothertown.''

\subsubsection{Convergence and Divergence}

Most Fielding (2006) entries with FF attestation can be mapped to cognates in
Prince-Speck (1904). The primary differences are systematic orthographic
conventions: Prince-Speck's \textit{g} initial corresponds to Fielding's
\textit{k} (voicing neutralization); Prince-Speck's \textit{ch} corresponds to
Fielding's \textit{c} (just a spelling difference); Prince-Speck's vowel
notation is inconsistent (reflecting English spelling conventions) while
Fielding's is phonemic; and Prince-Speck uses \texttt{'} for schwa while
Fielding omits schwa or encodes it as \textit{u}, \textit{a}, or \textit{i}.

\subsection{Current Orthographic Conventions}

\subsubsection{The Modern Practical Orthography}

The orthography used by Fielding (2006) and adopted by the Mohegan Language
Project is a \textbf{Latin-based phonemic alphabet} with the following
conventions:

\textbf{Alphabet (12 consonant symbols, 6 vowel symbols):}
\texttt{a â c e ê h i k m n ô p q s sh t u w y '}

Key design principles:
\begin{enumerate}
  \item \textbf{1:1 Phoneme-to-grapheme correspondence} --- each symbol maps
    to exactly one phoneme
  \item \textbf{Circumflex for length} --- â, ê, ô are long; a, e, i, u are
    short
  \item \textbf{\textit{c} for /tʃ/} --- avoids the English \textit{ch}
    digraph; never represents /k/
  \item \textbf{\textit{q} for /kʷ/} --- distinguishes from \textit{kw}
    cluster; occurs word-finally
  \item \textbf{\texttt{'} (okina) for glottal stop} --- follows Hawaiian and
    other indigenous orthography conventions
  \item \textbf{\textit{sh} for /ʃ/} --- standard English digraph
  \item \textbf{No silent letters} --- every written symbol is pronounced
  \item \textbf{Bound morphemes marked with \texttt{-}} --- prefixes:
    \textit{nu-}, \textit{ku-}; suffixes: \textit{-côq}
  \item \textbf{Part of speech tags} --- VAI, VII, VTA, VTI, NI, NA following
    Algonquianist conventions
  \item \textbf{Alternate spellings in parentheses} --- showing earlier
    English-based forms
\end{enumerate}

\subsubsection{Comparison with Granberry's Orthography}

Julian Granberry (2003) uses a somewhat different practical orthography,
differing primarily in representing lenis-fortis distinctions:

\begin{table}[h]
\centering
\begin{tabular}{lll}
\toprule
Fielding & Granberry & IPA \\
\midrule
p & b, p & /p/ \\
t & d, t & /t/ \\
k & g, k & /k/ \\
c & j, ch & /tʃ/ \\
s & z, s & /s/ \\
sh & zh, sh & /ʃ/ \\
\bottomrule
\end{tabular}
\caption{Fielding vs. Granberry orthography}
\label{tab:1364-granberry}
\end{table}

Granberry's orthography explicitly notes the vowel /ə/ (schwa) as \textit{u}
--- e.g., \textit{wutche} vs. Fielding's \textit{wicki}. He also uses
\textit{x} for /h/ in some environments. However, the Fielding orthography
has become the standard through the Mohegan Language Project's adoption.

\subsection{Problems Encountered}

\subsubsection{Small Corpus}

Only 145 records with explicit Mohegan-to-phoneme field mappings. This is the
smallest corpus among the three SNEA sources, limiting statistical approaches.

\subsubsection{Alternate Spellings}

The primary conversion challenge is English-based alternate spellings →
Fielding orthography, not Fielding → IPA. The alternate spellings provide
direct training data for converting Prince-Speck (1904) English-orthography
forms.

\subsubsection{Circularity Risk}

Some Fielding dictionary entries are reconstructed FROM Prince-Speck (1904)
--- using Fielding to verify Prince-Speck is circular for those entries.
Careful provenance tracking is required to distinguish entries with
independent FF attestation from those reconstructed via comparative methods.

\subsubsection{The Pequot Extinction and the Treaty of Hartford (1638)}

The Pequot War (1636--1638) culminated in the \textbf{Mystic Massacre} (May
26, 1637), where English forces and allies set fire to a Pequot fortified
village, blocking the exits and shooting anyone attempting to escape.
Approximately 700 Pequots were killed or captured; survivors were sold into
slavery in Bermuda and the West Indies, while others were dispersed as
captives to the victorious tribes.

The \textbf{Treaty of Hartford} (September 21, 1638) explicitly sought to
eradicate Pequot cultural identity:
\begin{itemize}
  \item \textbf{Name prohibition:} Survivors could no longer call themselves
    Pequots --- they were to be referred to as Mohegans or Narragansetts
    ``for ever''
  \item \textbf{Land prohibition:} Pequots were forbidden from returning to
    their native country
  \item \textbf{Dispersion:} Surviving Pequots were divided: ~80 to Mohegan
    Sachem Uncas, ~80 to Narragansett Sachem Miantonomo, ~20 to Niantic
    Sachem Ninigret
  \item \textbf{English control of Pequot country:} The Pequot lands went to
    the Connecticut River towns
\end{itemize}

The treaty effectively declared the Pequot nation extinct as a sovereign
polity. A 1743 copy of the treaty survives and is housed at the Connecticut
State Archives. Daragh Grant (2015, \emph{William and Mary Quarterly})
reconstructed the original 13-provision text from multiple surviving copies,
including a newly rediscovered 1665 copy in the British Library's Lansdowne
Manuscripts.

The Treaty of Hartford directly caused the language's near-extinction. By
outlawing Pequot speech and dispersing the population, the colonial government
ensured the language would cease to be transmitted to children. Mohegan
survived longer (until 1908) precisely because the Mohegan tribe was not
subject to the same dispersal --- their territory west of the Thames River
remained intact.

\subsection{Research Approach}

\subsubsection{Recommended Methodology}

Direct 1:1 phoneme mapping. Fielding's orthography is already phonemic ---
conversion is a direct lookup. The primary challenge is mapping alternate
English-based spellings to the standard Fielding form.

For the 145 records in our corpus, the practical orthography's transparency
means:
\begin{itemize}
  \item Phonemic conversion from the orthography is essentially 1:1 (98\%+
    regular)
  \item The primary challenge is not decoding Fielding orthography → phonemes,
    but mapping alternate spellings → Fielding orthography
  \item The alternate spellings provide direct training data for converting
    Prince-Speck (1904) English-orthography forms
  \item Vowel quality for /e/ (plain \textit{e}) is the single remaining
    ambiguity --- the phoneme ranges from /i/ to /e/ and its exact quality
    may need contextual disambiguation
\end{itemize}

\subsection{Conversion Workflow}
\texttt{[Pending implementation --- see Issue \#1364]}

\subsection{Linguist Feedback}
\texttt{[Template --- content pending review]}

\subsection{Related Work: Colonial Recorder Perception Problem}

The core challenge of this research card --- English-speaking recorders
misperceiving and mis-transcribing Algonquian phonemes --- is a known
phenomenon documented across Eastern North America. Two researchers' work is
directly relevant:

\textbf{Dr. Keith Cunningham} (Linguist, Nanticoke Indian Tribe). His
doctoral dissertation \emph{A Phonological Analysis of Nanticoke With
Practical Applications for Language Revitalization} (Georgetown University)
analyzes three 18th-century Nanticoke word lists (Heckewelder 1785) using the
same methodology: historical phonology from colonial recorders, with the same
challenges of multilingual informants and orthographic complexity. His 2024
Algonquian Conference presentation \emph{A Revised Phonological Analysis of
the Heckewelder Vocabulary of Nanticoke} demonstrates that the colonial
recorder perception problem extends beyond SNEA into the broader Eastern
Algonquian family. Cunningham's work on Nanticoke revitalization (co-authoring
a Nanticoke textbook) parallels the practical application of the Fielding
dictionary for Mohegan language reclamation.

\textbf{Dr. Craig Kopris} (Independent scholar). His paper \emph{Les sons
wyandots perçus par des oreilles étrangères} (Wyandot sounds perceived by
foreign ears, \emph{Recherches amérindiennes au Québec}, 2014) is the exact
framing of the Fielding/Prince-Speck problem: how European recorders
systematically misperceived and mis-transcribed indigenous phonemes. His
\emph{Wyandot Phonology: Recovering the Sound System of an Extinct Language}
applies the same recovery methodology to an unrelated language family
(Iroquoian), demonstrating that the foreign-ear perception problem is
universal across colonial documentation of North American languages. His
Dictionary of Wyandot project (NEH award FN-255576-17) shows the long-term
application of these methods to language revitalization.

The Mohegan-Pequot corpus is unique among the SNEA sources in having a
last-speaker corpus (Fidelia Fielding's diaries) alongside a modern practical
orthography (Fielding 2013). This positions it as the most phonologically
reliable source --- but the Prince-Speck (1904) glossary, with its
German-influenced transcription, reintroduces the foreign-ear perception
problem that Kopris documents for Wyandot.

\subsection{Mohegan Revitalization}

\subsubsection{The Mohegan Language Project}

Following the death of Fidelia Fielding in 1908, Mohegan had no fluent
speakers for nearly a century. In April 1998, the Council of Mohegan Tribal
Elders resolved to use Jits Bodunaxa's speech as the basis for a restored
Modern Mohegan (Granberry 2003:6). A Mohegan Language Committee was formed,
and Julian Granberry was appointed Tribal Linguist from 1998 to 2001.

In the 21st century, the \textbf{Mohegan Language Project} was established,
creating lessons, a dictionary, and online learning materials. The project
developed a practical orthography using the Latin alphabet, based on the four
phonetic diaries left by Fielding.

\subsubsection{Current Revitalization Efforts (as of 2024)}

The Mohegan Tribe operates an active \textbf{Language Restoration Project}
that ramped up significantly in 2017. Key initiatives include:
\begin{itemize}
  \item \textbf{Master-Apprentice Program:} Immersion-style learning with
    passive speaking → circumlocution → active speaking
  \item \textbf{Immersion Nests:} Yearlong cohort programs scaffolded to
    seasonal activities
  \item \textbf{Workbook Development:} First Mohegan language workbooks and
    grammar curricula
  \item \textbf{Two fluent teachers} (the most since Fidelia Fielding's time)
    leading classes
  \item \textbf{Online classes} open to all Mohegan tribal members and Mohegan
    household members
\end{itemize}

\subsubsection{Collaboration with Jessie Little Doe Baird}

MacArthur Fellow \textbf{Jessie Little Doe Baird} (Mashpee Wampanoag),
renowned for her Wôpanâak language reclamation, has served as a guide for the
Mohegan project. Using Wôpanâak (an N-dialect) as a sister language, the
Mohegan teachers cross-reference sound changes (N→Y correspondence) to build
vocabulary and grammar.

\subsubsection{The Samson Occom Papers}

In 2022, Dartmouth College repatriated the papers of \textbf{Samson Occom}
(1723--1792), a Mohegan minister and scholar who wrote extensively in his
native language. These 18th-century documents provide another primary source
layer for the Language Restoration Project, offering examples of the language
written by a fluent speaker during the period when Mohegan was still actively
spoken.

\subsubsection{Terminology: Reclamation vs. Revitalization}

Mohegan language teachers \textbf{Allyson Corbin} and \textbf{Autumn Cholewa},
who created much of the current curriculum, emphasize the term
\textbf{reclamation} over revitalization or revival:

\begin{quote}
``Reclaim is much more empowering for a community. The connotation is more
hands-on. More like, `This is up to you to do.' To reclaim it is to bring it
home. It's a prideful thing instead of something people felt pressure not to
do.'' (CT Insider, Nov 2024)
\end{quote}

\subsection{Bibliography}

\begin{enumerate}

\item Baird, Jessie Little Doe. (2010). MacArthur Fellowship for Wôpanâak
Language Reclamation.

\item Bloomfield, Leonard. (1925). ``On the Sound-System of Central
Algonquian.'' \emph{Language} 1:130--156.

\item Bloomfield, Leonard. (1946). ``Algonquian.'' In \emph{Linguistic
Structures of Native America}, ed. Harry Hoijer, 85--129. Viking Fund
Publications in Anthropology 6.

\item Costa, David. (2006). Lexical and grammatical research for the Modern
Mohegan Dictionary (unpublished ms).

\item Fielding, Stephanie. (2006). \emph{A Modern Mohegan Dictionary.}
Mohegan Tribe.

\item Fielding, Stephanie. (2013). \emph{Modern Mohegan Dictionary} (revised
edition). Mohegan Tribe.

\item Goddard, Ives. (1978). ``Eastern Algonquian Languages.'' In
\emph{Handbook of North American Indians}, vol. 15: Northeast, ed. Bruce G.
Trigger, 70--77. Smithsonian Institution.

\item Goddard, Ives. (1979). ``Comparative Algonquian.'' In \emph{The
Languages of Native America}, ed. Lyle Campbell and Marianne Mithun, 70--132.
University of Texas Press.

\item Granberry, Julian. (2003). \emph{Modern Mohegan: The Dialect of Jits
Bodunaxa.} Languages of the World/Materials 430. LINCOM Europa.

\item Granberry, Julian. (2003). \emph{A Lexicon of Modern Mohegan: The
Dialect of Jits Bodunaxa.} Languages of the World/Dictionaries 28. LINCOM
Europa.

\item Grant, Daragh. (2015). ``The Treaty of Hartford (1638): Reconsidering
Jurisdiction in Southern New England.'' \emph{William and Mary Quarterly}
72(3):461--498.

\item Haas, Mary. (1954). ``The Proto-Algonquian Word for `Sun.''' In
\emph{Language in Culture and Society}, ed. Dell Hymes. Harper \& Row.

\item Oxford, Will. (2015). ``Patterns of Contrast in Phonological Change:
Evidence from Algonquian Vowel Systems.'' \emph{Language} 91(2):308--339.

\item Oxford, Will. (2023). ``Consonant Inventories from Proto-Algonquian to
the Daughter Languages.'' University of Manitoba Algonquian Linguistics
Resources.

\item Pentland, David H. (1979). \emph{Algonquian Historical Phonology.} PhD
dissertation, University of Toronto.

\item Prince, J. Dyneley \& Speck, Frank G. (1903). ``The Mohegan-Pequot
Language.'' \emph{American Anthropologist} 5:193--212.

\item Prince, J. Dyneley \& Speck, Frank G. (1904). ``Glossary of the
Mohegan-Pequot Language.'' \emph{American Anthropologist} 6(1):18--44.

\item Siebert, Frank T. (1975). ``Resurrecting Virginia Algonquian from the
Dead.'' In \emph{Studies in Southeastern Indian Languages}, ed. James M.
Crawford, 285--453. University of Georgia Press.

\item Speck, Frank G. (1928). ``Native Tribes and Dialects of Connecticut.''
\emph{Annual Report of the Bureau of American Ethnology} 43:199--287.

\item Valentine, J. Randolph. (2001). \emph{Nishnaabemwin Reference
Grammar.} University of Toronto Press.

\end{enumerate}
